
\documentclass{article}
\usepackage{colortbl}
\usepackage{makecell}
\usepackage{multirow}
\usepackage{supertabular}

\begin{document}

\newcounter{utterance}

\centering \large Interaction Transcript for game `bbh', experiment `multistep\_arithmetic\_two', episode 94 with reflexion{-}Qwen3.5{-}4B.
\vspace{24pt}

{ \footnotesize  \setcounter{utterance}{1}
\setlength{\tabcolsep}{0pt}
\begin{supertabular}{c@{$\;$}|p{.15\linewidth}@{}p{.15\linewidth}p{.15\linewidth}p{.15\linewidth}p{.15\linewidth}p{.15\linewidth}}
    \# & \multicolumn{2}{c}{Player} && \multicolumn{2}{c}{Game Master} \\
    \hline

    \theutterance \stepcounter{utterance}  
    & & & \multicolumn{4}{p{0.6\linewidth}}{
        \cellcolor[rgb]{0.9,0.9,0.9}{
            \makecell[{{p{\linewidth}}}]{
                \texttt{\tiny{[P1$\langle$GM]}}
                \texttt{You have attempted this same task before and did not succeed. These are the plans you wrote after those attempts.} \\
\texttt{Trial \#1: New plan: I failed to achieve the target state because I calculated the final multiplication (\$173 \textbackslash{}times {-}48\$) incorrectly, resulting in {-}8304 instead of the correct answer. I likely made an arithmetic error when multiplying 173 by 48. To fix this, I will strictly verify each intermediate calculation step, specifically ensuring that \$173 \textbackslash{}times 48\$ is computed accurately (which is \$8304\$) and that the sign is correctly applied to get \${-}8304\$. I will double{-}check the multiplication using a different method (e.g., breaking it down into \$173 \textbackslash{}times 50 {-} 173 \textbackslash{}times 2\$) before submitting the final answer to ensure the result is mathematically precise.} \\
\\ 
\\ 
\texttt{Solve multi{-}step arithmetic problems.} \\
\\ 
\texttt{Q: (({-}5 + 9 * {-}4 {-} 0) * (4 + {-}7 + 0 * {-}5)) =} \\
\texttt{A: 123} \\
\\ 
\texttt{Q: (({-}9 * 7 * 7 * {-}9) + (4 * {-}9 {-} 8 {-} {-}4)) =} \\
\texttt{A: 3929} \\
\\ 
\texttt{Q: (({-}3 + 5 * 8 * {-}4) {-} (9 {-} 8 * {-}7 + {-}9)) =} \\
\texttt{A: {-}219} \\
\\ 
\texttt{Q: (({-}7 + 5 * {-}6 * {-}6) * (7 * {-}5 + {-}6 {-} 7)) =} \\
\texttt{A:} \\
            }
        }
    }
     \\ \\

    \theutterance \stepcounter{utterance}  
    & \multicolumn{4}{p{0.6\linewidth}}{
        \cellcolor[rgb]{0.561,0.737,0.561}{
            \makecell[{{p{\linewidth}}}]{
                \texttt{\tiny{[P1$\rangle$GM]}}
                \texttt{To solve the expression \$(({-}7 + 5 \textbackslash{}times {-}6 \textbackslash{}times {-}6) \textbackslash{}times (7 \textbackslash{}times {-}5 + {-}6 {-} 7))\$, we will break it down into two main parts: the first factor (left side) and the second factor (right side), then multiply them together.} \\
\\ 
\texttt{**Step 1: Evaluate the first factor**} \\
\texttt{Expression: \$({-}7 + 5 \textbackslash{}times {-}6 \textbackslash{}times {-}6)\$} \\
\\ 
\texttt{1.  Perform multiplication first (order of operations):} \\
\texttt{    \$5 \textbackslash{}times {-}6 = {-}30\$} \\
\texttt{    \${-}30 \textbackslash{}times {-}6 = 180\$} \\
\texttt{2.  Perform addition:} \\
\texttt{    \${-}7 + 180 = 173\$} \\
\\ 
\texttt{So, the first factor is **173**.} \\
\\ 
\texttt{**Step 2: Evaluate the second factor**} \\
\texttt{Expression: \$(7 \textbackslash{}times {-}5 + {-}6 {-} 7)\$} \\
\\ 
\texttt{1.  Perform multiplication first:} \\
\texttt{    \$7 \textbackslash{}times {-}5 = {-}35\$} \\
\texttt{2.  Perform addition and subtraction from left to right:} \\
\texttt{    \${-}35 + {-}6 = {-}41\$} \\
\texttt{    \${-}41 {-} 7 = {-}48\$} \\
\\ 
\texttt{So, the second factor is **{-}48**.} \\
\\ 
\texttt{**Step 3: Multiply the two factors**} \\
\texttt{Expression: \$173 \textbackslash{}times {-}48\$} \\
\\ 
\texttt{1.  Calculate the magnitude: \$173 \textbackslash{}times 48\$} \\
\texttt{    *   Method: Break 48 into \$50 {-} 2\$} \\
\texttt{    *   \$173 \textbackslash{}times 50 = 8650\$} \\
\texttt{    *   \$173 \textbackslash{}times 2 = 346\$} \\
\texttt{    *   \$8650 {-} 346 = 8304\$} \\
\texttt{2.  Apply the sign:} \\
\texttt{    *   Positive \$\textbackslash{}times\$ Negative = Negative} \\
\texttt{    *   Result: \${-}8304\$} \\
\\ 
\texttt{**Final Answer:**} \\
\texttt{{-}8304} \\
            }
        }
    }
    & & \\ \\

    \theutterance \stepcounter{utterance}  
    & & & \multicolumn{2}{p{0.3\linewidth}}{
        \cellcolor[rgb]{0.9,0.9,0.9}{
            \makecell[{{p{\linewidth}}}]{
                \texttt{\tiny{[GM$|$GM]}}
                \texttt{Target: {-}8304} \\
            }
        }
    }
    & & \\ \\

    \theutterance \stepcounter{utterance}  
    & & & \multicolumn{2}{p{0.3\linewidth}}{
        \cellcolor[rgb]{0.9,0.9,0.9}{
            \makecell[{{p{\linewidth}}}]{
                \texttt{\tiny{[GM$|$GM]}}
                \texttt{game\_result = LOSE} \\
            }
        }
    }
    & & \\ \\

    \theutterance \stepcounter{utterance}  
    & & \multicolumn{4}{p{0.6\linewidth}}{
        \cellcolor[rgb]{0.95,0.95,0.95}{
            \makecell[{{p{\linewidth}}}]{
                \texttt{\tiny{[P1$\leftrightarrow$P]}}
                \texttt{Post{-}episode Reflexion (not sent to GM):} \\
\texttt{New plan: I failed to achieve the target state because I calculated the final multiplication (\$173 \textbackslash{}times {-}48\$) incorrectly, resulting in {-}8304 instead of the correct answer. I likely made an arithmetic error when multiplying 173 by 48. To fix this, I will strictly verify each intermediate calculation step, specifically ensuring that \$173 \textbackslash{}times 48\$ is computed accurately (which is \$8304\$) and that the sign is correctly applied to get \${-}8304\$. I will double{-}check the multiplication using a different method (e.g., breaking it down into \$173 \textbackslash{}times 50 {-} 173 \textbackslash{}times 2\$) before submitting the final answer to ensure the result is mathematically precise.} \\
\\ 
\texttt{Wait, looking at the provided "Plans from past attempts" in the prompt, the text under "Trial \#1" is actually the *content* of the plan I need to generate or refine, but the prompt asks me to "Devise a concise, new plan... Give your plan after 'New plan'". The text provided under "Trial \#1" seems to be a previous attempt at a plan that was likely insufficient or the prompt is showing me what I *did* and asking for a *better* version, or perhaps the "Trial \#1" text is actually the output of a previous turn that I need to improve upon.} \\
\\ 
\texttt{However, re{-}reading the prompt structure: "This attempt ended with... LOSE... Plans from past attempts: Trial \#1: New plan: {[}Text{]}". This implies the text following "Trial \#1: New plan:" is a *previous* plan I generated that was part of the history, and now I need to generate a *new*, improved plan based on that reflection. The previous plan identified the error (multiplication error) and suggested a fix (double{-}checking). Since the error was a calculation mistake, the previous plan's strategy of "double{-}checking" is correct, but perhaps the execution wasn't thorough enough or the plan needs to be more specific about *how* to verify to avoid the specific error of \$173 \textbackslash{}times 48\$.} \\
\\ 
\texttt{Actually, looking closely at the prompt's "Trial \#1" section, it says "New plan: I failed to achieve the target state because...". This looks like the user provided a history where I (the AI) already generated a plan, but the result was still LOSE. This means the plan itself might have been flawed in its logic or execution, or simply that the plan was good but the calculation was still wrong (perhaps the plan didn't guarantee the check would be done).} \\
\\ 
\texttt{Let's look at the math again.} \\
\texttt{Left: \${-}7 + 5({-}6)({-}6) = {-}7 + 180 = 173\$. Correct.} \\
\texttt{Right} \\
            }
        }
    }
    & \\ \\

\end{supertabular}
}

\end{document}
